\section{Conclusion}\label{sec:conclusion}

Lattice Yang--Mills theory in four dimensions has a mass gap with the explicit
lower bound $m_G(\beta)$ for every compact simple gauge group in the range
$N\beta<\frac2{29}\kappa_G$, and for $\SU(2)$ up to $\beta_W=\frac13$ with
the bound $-\log(3\beta_W)$ and up to $\beta_W=\frac9{25}$ by resampling stars. A
single
spectral gap $\varpi$ of the dynamics that resamples cubes of a fixed physical
size $\ell$, uniform along the continuum tuning, gives every nontrivial
continuum limit the mass gap $16\varpi/(625e^2\ell)$. That spectral gap holds
at strong coupling, for $\SU(2)$ up to $\beta_W=\frac7{20}$ uniformly in the
cube size and for $\SU(3)$ up to $\beta_W=0.45$, and for massive free fields;
it fails for massless fields, at fixed cube size at weak coupling, and in the
abelian Coulomb phase.

The cube gap follows from a decay of correlations, below the sharp threshold
$\frac12$, between the gauge-invariant observables on the two sides of a slab
of physical width $\ell/2$, uniformly in the boundary condition. This
hypothesis is a form of mass generation at one scale. It holds in two
dimensions along the weak-coupling tuning and fails there for strips one link
wide; in four dimensions it is not proved. Curvature arguments, with or
without gauge fixing, stop at explicit couplings of order one; spectral gaps
for blocks of fixed lattice size force a trivial continuum limit; and a
lattice gap tending to zero at the right rate produces a gap in a continuum
limit that is not identified with Yang--Mills theory. Statements (D) and (E)
remain open, and with them the Clay problem; what this paper contributes is an
explicit strong-coupling theory for every gauge group, a reduction of the
continuum gap to one estimate on one scale and of that estimate to a decay of
correlations across one slab, and a map of what the present methods can reach.
